\documentclass[11pt,a4paper]{article}
\usepackage[margin=23mm]{geometry}
\usepackage{amsmath,amssymb,amsthm}
\setlength{\emergencystretch}{2em}
\newtheorem{lemma}{Lemma}
\newtheorem{theorem}[lemma]{Theorem}
\newcommand{\gen}[1]{\langle#1\rangle}
\title{The endpoint equality paths force an opposite head}
\author{Human Compression project: L044 continuation}
\date{}
\begin{document}\maketitle
Retain the complete critical-head hypotheses and acquired proofs,
including the explicit smaller-degree alternating induction. In
particular $n=2r$, the head is supported on $D$, and
\[
 R=G_{\ge3}=\operatorname{Alt}(\Omega\setminus(D\setminus\{6\})).
\]
The actual prefixes $P_i=G_{\le i}$ have their proved pure
alternating kernels and exact intersection formula with $R$.
Eligibility always concerns full ambient omissions.

\begin{lemma}[All later marks fix a saturated prefix support]
\label{fix}
Suppose $g\ge5$, $P=P_{g-1}$ has exactly two nontrivial
orbits $K=C_{g-1}$ and $E$, and
$|K|+|E|=2g+1$. Suppose every label $3,\ldots,g-1$ is paired.
Then every $\rho_k$ with $k\ge g+1$ fixes $K\cup E$ pointwise.
\end{lemma}
\begin{proof}
Such a mark $z=\rho_k$ centralizes $P$, fixes $K$ pointwise
by its pure natural alternating kernel, and preserves the only
other nontrivial $P$-orbit $E$. If its restriction $f$ to $E$
is nonidentity, centralization and transitivity make $f$ a free
involution. Write $|E|=2m$.

Set $I=\{3,\ldots,g-1\}$ and $H=G_I$. Its two nontrivial
orbits are $T=K\setminus(D\setminus\{6\})$ and $E$,
with $|T|=|K|-5$, by the exact prefix intersection formula.
Let $L\subseteq I$ be the labels having an ambient-eligible
edge in $E$. Such an edge cannot be a pair of $f$, since
$z\in G_i$ for every $i\in I$. Select one such edge per
label in $L$ and project to the $m$ pairs of $f$.
The selected graph is a forest: an avoiding path around a
cycle lifts to $E$ because $H$ commutes with $f$, with
$z|_E=f$ supplying any switches within a pair. It would give
an actual avoiding path for an eligible edge. Thus $|L|\le m-1$.
Every eligible edge for a label in $I\setminus L$ lies in $T$.
The acquired selected-localization theorem gives
$2|I\setminus L|\le |T|-1=|K|-6$. Consequently
\[
 g-3\le \frac{|E|}{2}-1+\frac{|K|-6}{2}
             =\frac{2g-7}{2}<g-3,
\]
a contradiction. Hence $f=1$, proving the lemma. This is the
exterior-deficit argument for an arbitrary later mark; no
crossing or parity assumption on $z|_E$ is needed.
\end{proof}

\begin{theorem}[The first later endpoint cut]
\label{endpoint}
Suppose $B=G_{\ge4}$ is intransitive on $\Omega\setminus D$,
and $r-1$ is the least label at least five having a unique
eligible edge. Suppose its first shared label is $g\ge5$,
and $P=P_{g-1}$ is intransitive on its support. Then
$g=r-3$, and the actual last three generators, in reversed
order, form the pure degree-six $A_5$ marking of source
Table 2(2). In particular the endpoint branch reduces to an
actual opposite head, rather than a projected isomorphism.
\end{theorem}
\begin{proof}
The accepted early-cut calculation and the present least-cut
hypothesis make every label $3,\ldots,r-2$ paired. Write
$C=C_{r-2}$ for the head-side omission orbit and $A$ for the
other orbit. All indices below $g$ act on $C$ alone: none
is shared, and a change to a disjoint supported side would
disconnect the noncommutation diagram. Thus $P$ fixes $A$.
Put $Y=\operatorname{supp}(P)$ and $h=r-1-g\ge1$.

Apply the complete endpoint support-deficit lemma after full
reversal. Its crossing is now index zero, its last shared
index is $h$, and its remaining tail is exactly $P$.
It gives $|\operatorname{Fix}(P)|\ge2h+1$, hence
$|Y|\le2g+1$. This is the endpoint theorem and its endpoint
path equality, not an interior equality interface.
If $p\ge2$ counts the nontrivial $P$-orbits, then
$H=P\cap R=G_{3,\ldots,g-1}$ has support $|Y|-5$ and
$p$ nontrivial orbits. The acquired paired interval bound gives
\[
 2g-6\le |Y|-5-p\le2g-4-p.
\]
Therefore $p=2$, $|Y|=2g+1$, and $P$ has head orbit $K$
and one exterior orbit $E$. Lemma~\ref{fix} applies.

Write $s=\rho_{r-1}$ and $u_i=\rho_{r-1-i}$ for
$1\le i\le h$. The exact endpoint equality gives disjoint
paths
\[
 a_0\mathrel{-_{u_1}}a_1\mathrel{-_{u_2}}\cdots
       \mathrel{-_{u_h}}a_h,
 \qquad
 b_0\mathrel{-_{u_1}}b_1\mathrel{-_{u_2}}\cdots
       \mathrel{-_{u_h}}b_h,
\]
where $s$ crosses $a_0,b_0$,
\[
 A=\{a_0,\ldots,a_h\},\quad
 \operatorname{Fix}(P)=A\cup\{b_0,\ldots,b_{h-1}\},
 \quad b_h\in Y.
\]
These statements describe the exact path vertices and fixed set;
no unproved classification of all graph edges is inserted.

By Lemma~\ref{fix}, $s$ fixes $Y$ pointwise. If $h\ge2$,
start from $b_h$ and commute $s$ past
$u_h,u_{h-1},\ldots,u_2$ to see that it fixes
$b_{h-1},\ldots,b_1$ as well. Every one of these commutations
is at distance at least two in the original marking.

For $h\ge3$, the acquired exhausted anchored-path lemma,
applied to $u_1,\ldots,u_h$ on their $h+1$ points $A$,
gives the full natural symmetric action
\[
             \gen{u_2,\ldots,u_h}|_A
                    =\operatorname{Sym}(A\setminus\{a_0\}).
\]
Its hypotheses hold because the endpoint anchor $a_0$ is
fixed by $u_2,\ldots,u_h$, the displayed path makes the action
transitive, and the last mark is nonidentity there.
The unique crossing makes $s$ preserve $A\setminus\{a_0\}$;
it centralizes the displayed natural $S_h$, whose centralizer
is trivial for $h\ge3$. Thus it fixes that set pointwise.
We have accounted for every point: $s=(a_0\ b_0)$, contrary
to evenness. For $h=1$, the same conclusion follows because
$A\setminus\{a_0\}$ is a singleton and $Y$ is fixed.

It remains that $h=2$, so $g=r-3$. Here evenness of $s$,
its fixed set $Y\cup\{b_1\}$ and its unique crossing give
\[
 s=(a_0\ b_0)(a_1\ a_2).
\]
The mark $u_1=\rho_{g+1}$ also fixes $Y$ by
Lemma~\ref{fix}. Its displayed path edges and the involution
condition on the remaining single point of $A$ give
\[
 u_1=(a_0\ a_1)(b_0\ b_1),\qquad
 u_2=(a_1\ a_2)(b_1\ b_2)\delta,
\]
where $u_2=\rho_g$ fixes $b_0$ by the unique crossing and
far commutation, and $\delta$ is an even involution supported
on $Y\setminus\{b_2\}$. These are exactly the accepted
six-point Table 2(2) core permutations, in reversed order,
with a possible disjoint exterior factor on the third mark.
Their precise degree-six $A_5$ identification is already paid.

The prefix $P_g$ is transitive on $Y\cup\{b_1\}$.
Indeed it preserves that set and fixes $b_0$, while adjoining
$u_1$ merely connects $b_0$ to $b_1$ on $C$; if any other
$P_g$-orbit in that set remained, $P_{r-2}$ would not be
transitive on $C$. Its pure alternating kernel therefore
contains $\delta$.

If $\delta\ne1$, the subgroup $Q=\gen{s,u_1,u_2}$ projects
onto the displayed perfect core $A_5$ and onto $\gen\delta$.
If it did not contain the pure exterior $\delta$, projection
to the core would be injective and give a nontrivial homomorphism
$A_5\longrightarrow C_2$. This is impossible. Thus
$\delta\in Q$. Actual IP then gives
\[
            \delta\in Q\cap P_g=\gen{u_2}.
\]
But $\delta$ fixes $b_1$ whereas $u_2$ moves it. Hence
$\delta=1$, a contradiction. The exterior factor is therefore
trivial, and the actual opposite head is exactly as claimed.
\end{proof}

\paragraph{Attribution and boundary.}
The endpoint paths (source Proposition 5.20), exhausted symmetric
path and exact perfect core are acquired proofs. The new step
fixes the old support under every later mark. All predecessors
remain paid, including the root's parity-fibre and global-middle
contributions. The intransitive $B$ and prefix hypotheses remain
explicit. Opposite-head exclusion, the unshared singleton endpoint
and complete L044 recovery remain separate. No new computation,
formal theorem, whole-original or question credit is claimed.
\end{document}
